% !TEX root = ../main.tex
%% Données utiles (constantes, tableau périodique, électronégativités) et repères du programme

\editionPage{donnees-utiles}{Données utiles}{Constantes, conversions et données de référence}{Valeurs arrondies, suffisantes pour les exercices · Constantes fondamentales d'après le CODATA}
\begin{blocEdition}{donnees-utiles}
\fontsize{10}{13.5}\selectfont
\editionEtape{Constantes physiques}
{\small\begin{tabularx}{\linewidth}{@{}Yl|Yl@{}}\toprule
Célérité de la lumière dans le vide & $c=\qty{3,00e8}{m.s^{-1}}$ & Constante d'Avogadro & $N_\mathrm{A}=\qty{6,02e23}{mol^{-1}}$ \\
Constante de Planck & $h=\qty{6,63e-34}{J.s}$ & Constante des gaz parfaits & $R=\qty{8,314}{J.mol^{-1}.K^{-1}}$ \\
Charge élémentaire & $e=\qty{1,60e-19}{C}$ & Constante de Faraday & $F=\qty{9,65e4}{C.mol^{-1}}$ \\
Masse de l'électron & $m_e=\qty{9,11e-31}{kg}$ & Constante de gravitation & $G=\qty{6,67e-11}{N.m^2.kg^{-2}}$ \\
Masses du proton et du neutron & $m_p\approx m_n=\qty{1,67e-27}{kg}$ & Permittivité du vide & $\varepsilon_0=\qty{8,85e-12}{F.m^{-1}}$ \\
Produit ionique de l'eau (\qty{25}{\degreeCelsius}) & $K_e=\num{1.0e-14}$ & Constante de Stefan-Boltzmann & $\sigma=\qty{5,67e-8}{W.m^{-2}.K^{-4}}$ \\
Pesanteur (France) & $g=\qty{9,81}{m.s^{-2}}$ & Intensité sonore de référence & $I_0=\qty{1,0e-12}{W.m^{-2}}$ \\\bottomrule
\end{tabularx}}
\editionEtape{Conversions et données courantes}
{\small\begin{tabularx}{\linewidth}{@{}Yl|Yl@{}}\toprule
Électronvolt & $\qty{1}{eV}=\qty{1,60e-19}{J}$ & Volume molaire (\qty{0}{\degreeCelsius} ; \qty{1,013}{bar}) & \qty{22,4}{L.mol^{-1}} \\
Kilowattheure & $\qty{1}{kW.h}=\qty{3,6e6}{J}$ & Volume molaire (\qty{20}{\degreeCelsius} ; \qty{1,013}{bar}) & \qty{24,0}{L.mol^{-1}} \\
Pression & $\qty{1}{bar}=\qty{1e5}{Pa}$ ; $\qty{1}{atm}=\qty{1,013e5}{Pa}$ & Eau : masse volumique & \qty{1,00e3}{kg.m^{-3}} \\
Température & $T(\unit{K})=\theta(\unit{\degreeCelsius})+\num{273,15}$ & Eau : capacité thermique massique & \qty{4,18e3}{J.kg^{-1}.K^{-1}} \\
Année & $\qty{1}{an}\approx\qty{3,16e7}{s}$ & Célérité du son dans l'air à \qty{20}{\degreeCelsius} & \qty{343}{m.s^{-1}} \\
Terre : masse, rayon & \qty{5,97e24}{kg} ; \qty{6371}{km} & Célérité du son dans l'eau & \qty{1,48e3}{m.s^{-1}} \\
Soleil : masse & \qty{1,99e30}{kg} & Distance Terre-Soleil & \qty{1,50e11}{m} \\\bottomrule
\end{tabularx}}
\editionEtape{Quelques couples acide-base à 25~\textdegree C}
{\small\begin{tabularx}{\linewidth}{@{}Yc|Yc@{}}\toprule
$\ce{H3O+}/\ce{H2O}$ & 0 & $\ce{CO2},\ce{H2O}/\ce{HCO3-}$ & 6,4 \\
acide lactique / ion lactate & 3,9 & $\ce{H2PO4-}/\ce{HPO4^{2-}}$ & 7,2 \\
acide méthanoïque / ion méthanoate & 3,8 & $\ce{HClO}/\ce{ClO-}$ & 7,5 \\
acide ascorbique / ion ascorbate & 4,1 & $\ce{NH4+}/\ce{NH3}$ & 9,2 \\
acide benzoïque / ion benzoate & 4,2 & $\ce{HCO3-}/\ce{CO3^{2-}}$ & 10,3 \\
acide éthanoïque / ion éthanoate & 4,8 & $\ce{H2O}/\ce{HO-}$ & 14 \\\bottomrule
\end{tabularx}}
\end{blocEdition}

\editionPage{donnees-utiles-2}{Tableau périodique}{Masses molaires et électronégativités}{Éléments des cinq premières périodes}
\begin{blocEdition}{donnees-utiles-2}
\fontsize{10}{13.5}\selectfont
\label{figedition-donnees-2}\schema{annexe-tableau-periodique}
\editionEtape{Autres éléments rencontrés dans le livre}
{\small\begin{tabularx}{\linewidth}{@{}YYYYY@{}}\toprule
Cs : $Z=55$ ; 132,9 & Ba : $Z=56$ ; 137,3 & Pt : $Z=78$ ; 195,1 & Au : $Z=79$ ; 197,0 & Hg : $Z=80$ ; 200,6 \\
Pb : $Z=82$ ; 207,2 & Rn : $Z=86$ ; (222) & Ra : $Z=88$ ; (226) & U : $Z=92$ ; 238,0 & Pu : $Z=94$ ; (244) \\\bottomrule
\end{tabularx}}\par\smallskip
{\small Masses molaires en \unit{g.mol^{-1}} ; une valeur entre parenthèses est le nombre de masse de l'isotope le plus stable d'un
élément sans isotope stable.}
\editionEtape{Utiliser l'électronégativité}
L'électronégativité $\chi$ traduit la capacité d'un atome à attirer vers lui les électrons d'une liaison covalente. Dans une
liaison entre deux atomes d'électronégativités différentes, l'atome le plus électronégatif porte une charge partielle
négative $\delta^-$, l'autre une charge partielle positive $\delta^+$ : la liaison est polarisée. On considère en général qu'une
liaison est polarisée lorsque l'écart d'électronégativité dépasse 0,4 environ ; c'est le cas des liaisons C--O, O--H, C--N
ou C--Cl, mais pas de la liaison C--H. Les sites riches en électrons (atomes porteurs de $\delta^-$ ou de doublets non liants)
sont des sites \textbf{donneurs} de doublet ; les sites pauvres en électrons ($\delta^+$, charge positive) sont des sites
\textbf{accepteurs}.
\editionEtape{Spectroscopie infrarouge : nombres d'onde caractéristiques}
{\small\begin{tabularx}{\linewidth}{@{}Yl|Yl@{}}\toprule
O--H d'un alcool (liaison H) & 3200 -- \qty{3400}{cm^{-1}}, large & C$=$O d'un ester & 1735 -- \qty{1750}{cm^{-1}} \\
O--H d'un acide carboxylique & 2500 -- \qty{3200}{cm^{-1}}, très large & C$=$O d'un acide & 1700 -- \qty{1725}{cm^{-1}} \\
N--H d'une amine ou d'un amide & 3100 -- \qty{3500}{cm^{-1}} & C$=$O d'un aldéhyde, d'une cétone & 1650 -- \qty{1730}{cm^{-1}} \\
C--H & 2800 -- \qty{3100}{cm^{-1}} & C$=$O d'un amide & 1650 -- \qty{1700}{cm^{-1}} \\\bottomrule
\end{tabularx}}
\end{blocEdition}

\editionPage{couverture-programme}{Repères du programme}{Notions et capacités exigibles}{Programme de spécialité de terminale (BO spécial n°~8 du 25 juillet 2019) · Renvois aux chapitres}
\begin{blocEdition}{couverture-programme}
\fontsize{10}{13.5}\selectfont
Le programme est organisé en quatre thèmes, complétés par une partie transversale, « Mesure et incertitudes ». Les
tableaux suivent son plan et indiquent les chapitres du livre où chaque notion est traitée. Les capacités exigibles
détaillées sont reprises en tête de chaque chapitre, dans la rubrique « Objectifs ».
\editionEtape{Constitution et transformations de la matière}
{\small\begin{tabularx}{\linewidth}{@{}p{50mm}Yr@{}}\toprule
\textbf{Partie du programme} & \textbf{Notions principales} & \textbf{Ch.} \\\midrule
1.A Transformations acide-base & couple acide-base, réaction acide-base, pH, acides et bases forts & 1 \\
1.B Méthodes physiques d'analyse & loi de Beer-Lambert, loi de Kohlrausch, gaz parfait, spectroscopies IR et UV-visible & 2 \\
1.C Méthodes chimiques d'analyse & titrages avec suivi pH-métrique ou conductimétrique, équivalence & 3 \\
2.A Évolution temporelle & vitesses volumiques, temps de demi-réaction, facteurs cinétiques, catalyse, loi d'ordre 1,
mécanisme réactionnel, flèches courbes & 4, 5 \\
2.B Transformation nucléaire & stabilité des noyaux, désintégrations, loi de décroissance, datation, radioprotection & 6 \\
3.A Sens d'évolution spontanée & taux d'avancement final, quotient de réaction, constante d'équilibre, critère d'évolution,
piles & 7, 8 \\
3.B Force des acides et des bases & $K_\mathrm{A}$, produit ionique de l'eau, diagrammes de prédominance et de distribution,
solution tampon & 9 \\
3.C Forcer le sens d'évolution & électrolyse, quantité d'électricité, accumulateurs, stockage de l'énergie & 10 \\
4. Stratégies en synthèse organique & structure et nomenclature, isomérie, polymères, optimisation d'une synthèse,
protection de groupes, chimie verte & 11, 12 \\\bottomrule
\end{tabularx}}
\editionEtape{Mouvement et interactions}
{\small\begin{tabularx}{\linewidth}{@{}p{50mm}Yr@{}}\toprule
1. Décrire un mouvement & vecteurs position, vitesse et accélération ; repère de Frenet ; mouvements rectilignes et
circulaires & 13 \\
2. Relier les actions au mouvement & deuxième loi de Newton, centre de masse, référentiel galiléen & 14 \\
& mouvement dans un champ uniforme (pesanteur, champ électrique), aspects énergétiques & 15 \\
& mouvement dans un champ de gravitation, satellites, lois de Kepler & 16 \\
3. Modéliser l'écoulement d'un fluide & poussée d'Archimède, débit volumique, relation de Bernoulli, effet Venturi & 17 \\\bottomrule
\end{tabularx}}
\end{blocEdition}

\editionPage{couverture-programme-2}{Repères du programme (suite)}{Notions et capacités exigibles}{Programme de spécialité de terminale · Renvois aux chapitres}
\begin{blocEdition}{couverture-programme-2}
\fontsize{10}{13.5}\selectfont
\editionEtape{L'énergie : conversions et transferts}
{\small\begin{tabularx}{\linewidth}{@{}p{50mm}Yr@{}}\toprule
1. Décrire un système thermodynamique & échelles microscopique et macroscopique, modèle du gaz parfait, équation d'état,
limites du modèle & 18 \\
2. Bilans d'énergie, premier principe & énergie interne, travail, transfert thermique, capacité thermique & 19 \\
& modes de transfert, flux thermique, résistance thermique, bilan radiatif de la Terre, loi de Newton & 20 \\\bottomrule
\end{tabularx}}
\editionEtape{Ondes et signaux}
{\small\begin{tabularx}{\linewidth}{@{}p{50mm}Yr@{}}\toprule
1. Caractériser les phénomènes ondulatoires & intensité sonore, niveau d'intensité sonore, atténuations & 21 \\
& diffraction, interférences, fentes d'Young & 22 \\
& effet Doppler & 23 \\
2.A Former des images & lunette astronomique afocale, grossissement & 24 \\
2.B Décrire la lumière par un flux de photons & énergie du photon, effet photoélectrique, cellules photovoltaïques,
rendement & 25 \\
3. Dynamique d'un système électrique & intensité et débit de charges, condensateur, circuit RC, capteurs capacitifs & 26 \\\bottomrule
\end{tabularx}}
\editionEtape{Mesure et incertitudes}
{\small\begin{tabularx}{\linewidth}{@{}p{50mm}Yr@{}}\toprule
Partie transversale & variabilité de la mesure, histogramme, incertitudes-types de type A et B, incertitudes composées,
écriture du résultat, comparaison à une valeur de référence & 27 \\\bottomrule
\end{tabularx}}
\editionEtape{Capacités numériques du programme}
{\small\begin{tabularx}{\linewidth}{@{}Yr@{}}\toprule
Représenter l'histogramme associé à une série de mesures & 27 \\
Simuler un processus aléatoire illustrant la détermination d'une grandeur avec incertitudes-types composées & 27 \\
Représenter l'évolution des quantités de matière des espèces au cours d'un titrage & 3 \\
Tracer l'évolution temporelle d'une concentration, d'une vitesse volumique, et tester une loi d'ordre 1 & 4 \\
Déterminer le taux d'avancement final d'une transformation modélisée par la réaction d'un acide sur l'eau & 7, 9 \\
Tracer le diagramme de distribution des espèces d'un couple acide-base & 9 \\
Représenter des vecteurs accélération d'un point lors d'un mouvement & 13 \\
Représenter l'évolution des grandeurs énergétiques d'un système en mouvement dans un champ uniforme & 15 \\
Exploiter des données astronomiques ou satellitaires pour tester les lois de Kepler & 16 \\
Représenter la somme de deux signaux sinusoïdaux synchrones en faisant varier leur déphasage & 22 \\\bottomrule
\end{tabularx}}\par\smallskip
Le livre propose aussi des modélisations par la méthode d'Euler (chapitres 14, 20 et 26) et d'autres programmes courts dans
chaque chapitre : chacun est repéré par l'étiquette « Capacité numérique ».
\end{blocEdition}
